\section{Diseño y desarrollo del producto}
~\label{sec:desarrollo}

La presente sección describe el diseño y la implementación del sistema 
a partir de los requerimientos especificados, la arquitectura base 
definida durante la fase de Elaboración, y las decisiones técnicas 
adoptadas durante la fase de Construcción. La organización del 
contenido sigue la descomposición natural del sistema: primero el 
análisis de requerimientos (Sección~\ref{subsec:requerimientos}), 
luego el diseño arquitectónico global (Sección 
~\ref{subsec:arquitectura}), y a continuación la implementación de 
cada componente en el orden en que se integran durante el ciclo de 
vida del dato: el sensor de captura (Sección~\ref{subsec:sensor}), el 
backend y su pipeline de análisis (Sección~\ref{subsec:backend}), el 
motor de detección de anomalías (Sección~\ref{subsec:motor}), la capa 
de exposición mediante API REST y WebSocket (Sección 
~\ref{subsec:api}), y la infraestructura de integración y entrega 
continua (Sección~\ref{subsec:cicd}).

\subsection{Análisis y especificación de requerimientos}
~\label{subsec:requerimientos}

La especificación de requerimientos del sistema se documentó 
siguiendo el estándar IEEE 830, con identificadores únicos, 
prioridad de desarrollo, entradas y salidas, precondiciones, 
postcondiciones y flujos básico y alternativo para cada requerimiento 
funcional. El trabajo produjo 54 requerimientos funcionales (RF-01 a 
RF-54) distribuidos en los cuatro módulos del sistema, y seis 
categorías de requerimientos no funcionales que abarcan rendimiento, 
seguridad, fiabilidad, disponibilidad, mantenibilidad y portabilidad.

La Tabla~\ref{tab:rf-resumen} presenta la distribución de 
requerimientos por módulo junto con una síntesis de las 
responsabilidades cubiertas. La agrupación coincide con la 
descomposición arquitectónica descrita en la Sección 
~\ref{subsec:arquitectura} y permite trazar cada módulo a su conjunto 
de requerimientos asociados.

\begin{table}[ht]
\centering
\caption{Distribución de requerimientos funcionales por módulo del sistema.}
~\label{tab:rf-resumen}
\begin{tabular}{@{}p{2.3cm}cp{4cm}@{}}
\toprule
\textbf{Módulo} & \textbf{RF} & \textbf{Responsabilidades cubiertas} \\
\midrule
Sensor de captura & RF-01 a RF-05 & Captura de tráfico, métricas del endpoint, descubrimiento ARP, registro y consentimiento. \\[4pt]
Backend y pipeline & RF-06 a RF-25 & Recepción y persistencia, análisis de flujos, cálculo de métricas contextualizadas, inventario de activos. \\[4pt]
Motor de anomalías & RF-26 a RF-35 & Baselines por dispositivo, detección EWMA, filtro de persistencia, generación de alertas. \\[4pt]
API y WebSocket & RF-36 a RF-45 & Autenticación, endpoints de consulta, canal de eventos en tiempo real, notificaciones. \\[4pt]
Aplicación móvil & RF-46 a RF-54 & Login, dashboard, inventario, detalle de dispositivo, alertas, notificaciones push. \\
\bottomrule
\end{tabular}
\end{table}

Los requerimientos se priorizaron en tres niveles: \textit{Alta} para 
aquellos cuya ausencia impide la operación básica del sistema, 
\textit{Media} para los que amplían la utilidad sin ser 
bloqueantes, y \textit{Baja} para los que representan mejoras 
incrementales. Esta priorización orientó el orden de las iteraciones 
de Construcción y permitió organizar la batería de pruebas para 
cubrir primero las rutas críticas del sistema.

\subsection{Diseño arquitectónico}
~\label{subsec:arquitectura}

La arquitectura del sistema se describe en tres niveles de 
abstracción, siguiendo el modelo C4 de Simon Brown: contenedores, 
componentes y despliegue. El nivel de contexto (actores externos) se 
aborda implícitamente en la Sección~\ref{sec:alcance}, donde se 
caracterizaron el administrador como único actor humano y los 
endpoints como elementos de entrada al sistema.

\subsubsection{Vista de contenedores y componentes}

El sistema se descompone en tres subsistemas principales que se 
comunican mediante canales bien definidos. El primero es el sensor 
de captura, desplegado en cada endpoint monitoreado, que expone dos 
interfaces hacia el backend: \texttt{IEndpointMetrics} para el envío 
periódico de métricas del sistema operativo por UDP, e 
\texttt{IPcapStream} para la transmisión continua de paquetes 
sanitizados por TCP.@ El segundo es el backend centralizado, que 
internamente se divide en cinco subsistemas de propósito específico 
conectados mediante interfaces internas. El tercero es la aplicación 
móvil, que consume las interfaces públicas del backend 
(\texttt{IRestApi} e \texttt{IWebSocket}) para presentar información 
al administrador.

\begin{figure}[H]
\centering
\includegraphics[width=\columnwidth]{diagrama_componentes.png}
\caption{Vista de componentes del sistema (modelo C4 nivel 3). Se muestran los cinco subsistemas internos del backend, las interfaces que exponen y las dependencias entre ellos.}
~\label{fig:componentes}
\end{figure}

La Fig.~\ref{fig:componentes} ilustra la estructura interna del 
backend. El subsistema de recepción (\texttt{UDP Receiver} y 
\texttt{TCP Receiver}) recibe los flujos provenientes de los 
sensores y los entrega a los subsistemas de análisis y almacenamiento 
intermedio. El subsistema de \textit{intake} rota los archivos PCAP 
en ventanas temporales manejables y los pone a disposición del motor 
de análisis mediante la interfaz \texttt{IRotatedPcap}. El subsistema 
de análisis combina la sonda ISP y el motor de análisis principal, 
que procesa los PCAP, genera las métricas contextualizadas y las 
entrega al subsistema de gestión a través de la interfaz 
\texttt{IProcessedData}. El subsistema de gestión consolida el 
inventario de activos y la generación de alertas. Finalmente, el 
subsistema de API, que agrupa el \texttt{Auth Manager}, el 
\texttt{API Handler} y el \texttt{WebSocket Handler}, expone los 
datos persistidos hacia la aplicación móvil. Todos los subsistemas 
acceden a la base de datos PostgreSQL mediante la interfaz 
\texttt{IRepository} expuesta por el subsistema \texttt{dbmanager}.

Esta descomposición responde a dos decisiones arquitectónicas clave. 
La primera es la separación del pipeline de procesamiento y la capa 
de exposición como procesos independientes del sistema operativo, 
comunicados exclusivamente a través de la base de datos compartida 
~\cite{hohpe2003eip}. Esto permite que un fallo en el pipeline no 
afecte la disponibilidad de la API para consultas móviles, y que 
ambos procesos escalen sus recursos de forma independiente según su 
perfil de carga~\cite{silberschatz2018os}. La segunda es el uso de 
interfaces nombradas en el diagrama de componentes, que formalizan 
los contratos entre subsistemas y permiten sustituir implementaciones 
sin afectar al resto del sistema.

\subsubsection{Vista de despliegue}

\begin{figure*}[!t]
\centering
\includegraphics[width=0.95\textwidth]{diagrama_despliegue.png}
\caption{Vista de despliegue del sistema. Los sensores transmiten a un servidor central contenedorizado que expone servicios al dispositivo móvil mediante una red privada Tailscale y al servicio de notificaciones push Firebase Cloud Messaging.}
~\label{fig:diagrama_despliegue}
\end{figure*}

La Fig.~\ref{fig:diagrama_despliegue} presenta la disposición física del 
sistema en su entorno de operación. Los sensores de captura se 
ejecutan como artefactos \texttt{agente\_captura.py} sobre un runtime 
Python 3.x en cada uno de los endpoints de la red local. El servidor 
central, ubicado en la dirección estática 192.168.1.10, aloja tres 
contenedores Docker que ejecutan de forma independiente: el motor de 
pipeline (\texttt{goatguard\_motor.py}) en los puertos 9000 y 9001, 
la API REST (\texttt{api\_goatguard.py}) en el puerto 8000, y el 
servicio PostgreSQL 15 en el puerto 5432. La aplicación móvil, 
empaquetada como \texttt{app\_goatguard.apk}, se ejecuta sobre 
Android 10 o superior en el dispositivo del administrador.



Dos decisiones de despliegue merecen mención explícita. La primera 
es el uso de Tailscale como red privada virtual para conectar el 
dispositivo móvil con el servidor central, lo cual permite al 
administrador consultar el sistema desde fuera de la red local sin 
exponer el servidor a Internet mediante una IP pública. La segunda 
es la integración con Firebase Cloud Messaging para la entrega de 
notificaciones push al dispositivo móvil cuando el motor de análisis 
genera alertas que superan los umbrales de severidad.

\subsubsection{Modelo de datos}

La persistencia del sistema se organiza en torno a tres tipos 
estructurales de entidades que responden a perfiles de acceso 
distintos. Las entidades \textit{base} describen elementos 
relativamente estáticos del dominio cuya información cambia con poca 
frecuencia: una red, un dispositivo, un sensor registrado, un 
usuario autenticado. Las entidades \textit{acumuladas} conservan el 
histórico completo de las mediciones producidas por el motor de 
análisis a lo largo del tiempo; cada ciclo del motor inserta una 
nueva fila por métrica observada, y estas tablas crecen de forma 
monótona hasta que un proceso de retención las recorta. Las 
entidades \textit{sobreescritas} mantienen únicamente el valor más 
reciente de cada métrica y se actualizan mediante operaciones 
\texttt{UPSERT} que sustituyen la fila anterior por la nueva. La 
Fig.~\ref{fig:er-tipologia} ilustra los tres tipos sobre las tablas 
relacionadas con el monitoreo de una red.

\begin{figure}[H]
\centering
\includegraphics[width=\columnwidth]{er_tipologia.png}
\caption{Tipología de entidades del modelo de datos: \textsc{Network} como entidad base, \textsc{Network\_Snapshot} como entidad acumulada que preserva el histórico, y \textsc{Network\_Current\_Metrics} como entidad sobreescrita que almacena únicamente la última observación.}
~\label{fig:er-tipologia}
\end{figure}

Esta distinción responde a un compromiso entre la necesidad de 
mantener un historial para el análisis retrospectivo y la 
eficiencia de las consultas en tiempo real desde la aplicación 
móvil. Las entidades acumuladas operan como tablas de series 
temporales con inserciones continuas y lecturas analíticas 
ocasionales; su tamaño crece linealmente con el tiempo y requiere 
particionamiento y políticas de retención para mantenerse operativa. 
Las entidades sobreescritas, en contraste, mantienen una sola fila 
por dispositivo y se actualizan en cada ciclo, lo que permite al 
cliente móvil consultar el estado actual de la red en tiempo 
constante sin recorrer el histórico completo. El mismo patrón se 
aplica a los dispositivos (\textsc{Device\_Current\_Metrics} vs 
\textsc{Endpoint\_Snapshot}) y a los top talkers 
(\textsc{Top\_Talker\_Current} vs \textsc{Top\_Talker}).

El conjunto completo de entidades se agrupa en cuatro dominios 
funcionales, como se aprecia en el diagrama completo del modelo 
entidad-relación (Fig.~\ref{fig:er}, ubicada al final del 
documento). El dominio de inventario comprende \textsc{Network}, 
\textsc{Device} y \textsc{Agent}, que describen la estructura 
estática de la red monitoreada. El dominio de mediciones alberga 
las parejas acumulada y sobreescrita de cada tipo de métrica, junto 
con la entidad \textsc{ML\_Prediction} que almacena las 
clasificaciones producidas por el fingerprinting de tráfico. El 
dominio de análisis registra la salida del motor de anomalías: 
\textsc{Alert} para los eventos que superan los umbrales de 
severidad, \textsc{Insight} para las observaciones contextuales 
generadas durante el procesamiento, y \textsc{Top\_Talker} para la 
identificación de los dispositivos de mayor consumo. El dominio de 
autenticación agrupa \textsc{User}, \textsc{Session} y 
\textsc{Push\_Token}, responsables respectivamente de las 
credenciales, los tokens JWT activos y los identificadores para el 
envío de notificaciones push al dispositivo móvil del administrador.

\begin{figure*}[p]
\centering
\includegraphics[width=0.95\textwidth]{diagrama_er.png}
\caption{Modelo entidad-relación completo del sistema. Las tablas se agrupan en cuatro dominios funcionales: inventario, mediciones, análisis y autenticación.}
~\label{fig:er}
\end{figure*}

\subsection{Sensor de captura distribuida}
~\label{subsec:sensor}

El sensor de captura es el componente desplegado en cada endpoint 
de la red que observa el tráfico local, recolecta métricas del 
sistema y transmite ambos flujos hacia el backend central. Su 
diseño interno resuelve tres problemas técnicos de forma integrada: 
la captura continua sin bloquear la ejecución de tareas periódicas, 
la sanitización del tráfico antes de su transmisión para preservar 
la confidencialidad de las comunicaciones, y la entrega ordenada de 
paquetes al colector sobre un canal TCP que no preserva fronteras 
entre mensajes.

\subsubsection{Arquitectura interna}

La estructura del sensor (Fig.~\ref{fig:sensor-interno}) se 
organiza en cuatro módulos funcionales: un capturador de paquetes 
que intercepta el tráfico de la interfaz de red seleccionada, un 
sanitizador que aplica el recorte dinámico al payload y preserva la 
longitud original del paquete, un socket TCP que transmite los 
paquetes binarios sanitizados al colector por el puerto 9000, y un 
socket UDP que envía las métricas del endpoint (CPU, memoria RAM, 
disco, velocidad del enlace) serializadas en formato JSON 
~\cite{rfc8259} al puerto 9001.

\begin{figure}[H]
\centering
\includegraphics[width=\columnwidth]{sensor_interno.png}
\caption{Arquitectura interna del sensor de captura. El módulo de captura intercepta paquetes, el sanitizador los recorta, y dos sockets independientes transmiten tráfico binario (TCP) y métricas JSON (UDP) al colector.}
~\label{fig:sensor-interno}
\end{figure}

La selección de la interfaz de red principal se realiza de forma 
automática mediante un filtro que descarta las interfaces sin 
conectividad real: las de loopback, las que operan en el rango 
APIPA~\cite{rfc3927} (169.254.0.0/16, indicador de ausencia de 
asignación DHCP), las virtuales de hipervisores y contenedores, y 
los adaptadores Bluetooth. La interfaz seleccionada es aquella que 
posee simultáneamente una dirección IPv4 asignada y una dirección 
MAC asociada, criterio que evidencia conectividad activa a la 
infraestructura de red. El sensor construye su identificador único 
concatenando el hostname del sistema con la dirección MAC de esa 
interfaz, lo que garantiza unicidad a nivel de hardware y legibilidad 
humana en el inventario centralizado.

\subsubsection{Captura concurrente y modelo productor-consumidor}

El sensor ejecuta tareas con perfiles temporales distintos: la 
captura de paquetes es continua y bloqueante, mientras que el 
reporte de métricas y la señal de heartbeat operan en intervalos 
periódicos. La coordinación de estas tareas adopta un modelo de 
concurrencia basado en hilos, siguiendo el patrón productor-consumidor 
descrito por Buschmann et al.~\cite{buschmann1996posa}. Un hilo 
dedicado actúa como productor, invocando la función de captura sobre 
la interfaz seleccionada y depositando cada paquete procesado en un 
buffer intermedio. El hilo principal opera como consumidor, vaciando 
el buffer periódicamente para transmitir los paquetes acumulados al 
colector. La sincronización se implementa mediante exclusión mutua 
sobre el buffer compartido, lo que previene condiciones de carrera 
en el acceso concurrente~\cite{silberschatz2018os}.

El buffer se implementa como una cola acotada con política de 
descarte FIFO.@Cuando la capacidad máxima se alcanza, los paquetes 
más antiguos se descartan para acomodar los nuevos. Esta decisión 
protege la disponibilidad del endpoint frente a escenarios donde el 
consumidor no puede vaciar el buffer a la velocidad necesaria, como 
ocurre ante una pérdida temporal de conectividad con el colector: 
el consumo de memoria del sensor permanece acotado y predecible, en 
cumplimiento del requerimiento no funcional de rendimiento que 
establece un consumo máximo de memoria por sensor.

\subsubsection{Sanitización del tráfico}

El tráfico capturado no se transmite en su forma original al 
colector. Antes de su almacenamiento en el buffer, cada paquete 
atraviesa el sanitizador, que aplica un recorte dinámico al payload 
en función del puerto de destino del flujo. Los encabezados de las 
capas de enlace, red y transporte (aproximadamente 54 bytes, según 
se justificó en la Sección~\ref{sec:marco-teorico}) contienen toda 
la información necesaria para reconstruir flujos y extraer métricas 
de comportamiento, de modo que un recorte a 96 bytes preserva la 
totalidad de las cabeceras con margen para opciones TCP 
~\cite{stevens2004unix}. Para tráfico dirigido a los puertos 
asociados a DNS (53) y HTTPS (443), el recorte se extiende a 300 
bytes; este margen adicional captura el dominio consultado en el 
caso de DNS~\cite{rfc1035} y la indicación de nombre de servidor 
(SNI) del handshake TLS en el caso de HTTPS~\cite{rfc8446}, 
elementos con valor diagnóstico que no comprometen el contenido de 
las comunicaciones.

Cada paquete sanitizado conserva un campo \texttt{orig\_len} con la 
longitud original antes del truncamiento. Este campo es necesario 
para que el motor de análisis calcule métricas volumétricas como 
throughput, bytes por flujo, tamaño promedio de paquete y carga del 
enlace sobre el volumen real de tráfico, no sobre el volumen 
truncado. La sanitización materializa así el principio de 
minimización de datos exigido por el marco legal descrito en la 
Sección~\ref{sec:marco-legal}: se preserva exclusivamente la 
información técnica necesaria para el análisis, eliminando el 
contenido de las comunicaciones que podría constituir interceptación 
indebida.

\subsubsection{Protocolo binario de transporte}

TCP, como protocolo orientado a flujo de bytes, no preserva las 
fronteras entre los mensajes individuales emitidos por el transmisor 
~\cite{rfc793}. Un envío de 116 bytes por parte del sensor puede ser 
recibido por el colector como una lectura de 80 bytes seguida de 
otra de 36, o como una única lectura de 116, dependiendo del estado 
del buffer del kernel y de las condiciones de la red. Esta 
característica, conocida como \textit{message boundary problem} 
~\cite{stevens2004unix}, requiere que el protocolo de aplicación 
implemente su propio mecanismo de delimitación (Fig. 
~\ref{fig:message-boundary}).

\begin{figure}[H]
\centering
\includegraphics[width=\columnwidth]{tcp_framing.jpg}
\caption{Problema de fronteras en TCP y solución mediante \textit{length-prefix framing}. El receptor no puede inferir los límites de los mensajes originales sin un esquema explícito de delimitación. El prefijo de longitud antepuesto a cada payload resuelve la ambigüedad.}
~\label{fig:message-boundary}
\end{figure}

El sensor resuelve este problema mediante un esquema de 
\textit{length-prefix framing}: antepone a cada paquete capturado un 
encabezado fijo de 20 bytes que permite al receptor delimitar los 
mensajes dentro del flujo continuo de TCP.@El encabezado se 
estructura en \textit{network byte order} (big endian, conforme a 
RFC 793) y contiene cuatro campos (Fig.~\ref{fig:tcp-header}): la 
longitud original del paquete antes de la sanitización 
(\texttt{orig\_len}, 4 bytes), el puerto destino extraído de la 
cabecera de transporte del paquete capturado (\texttt{dst\_port}, 4 
bytes), la marca temporal de captura en formato IEEE 754 de doble 
precisión (\texttt{timestamp}, 8 bytes), y la longitud de los datos 
truncados que siguen al encabezado (\texttt{data\_len}, 4 bytes).

\begin{figure}[H]
\centering
\includegraphics[width=\columnwidth]{tcp_header.png}
\caption{Estructura del encabezado binario de transporte sensor-colector. Los 20 bytes de cabecera preceden a un payload de tamaño variable indicado por \texttt{data\_len}.}
~\label{fig:tcp-header}
\end{figure}

\subsubsection{Transporte dual y fundamentación de protocolos}

La separación del tráfico capturado (TCP) y las métricas del sistema 
(UDP) en dos canales paralelos responde a la diferencia en los 
requisitos de confiabilidad de cada tipo de dato. Los paquetes de red 
capturados son irrecuperables: cada paquete representa un evento 
único en la red, y si se pierde en tránsito no existe otra copia ni 
se regenera en el siguiente ciclo. La garantía de entrega ordenada y 
confiable que TCP proporciona mediante sus mecanismos de 
acknowledgement, retransmisión y control de flujo~\cite{rfc793} es 
por tanto apropiada para este canal. Las métricas del endpoint, en 
contraste, son tolerantes a pérdida: se generan periódicamente cada 
cinco segundos, cada paquete contiene una muestra completa del 
estado actual del sistema, y si un paquete UDP~\cite{rfc768} se 
pierde, el siguiente lo reemplaza con datos más recientes. La 
retransmisión de una métrica obsoleta consumiría ancho de banda sin 
aportar valor. El sobrecosto de ocho bytes de encabezado UDP frente 
a los veinte mínimos de TCP~\cite{stevens2004unix}, sumado a la 
ausencia de establecimiento de sesión, hace de UDP la elección 
natural para este perfil de transmisión.

\subsection{Backend centralizado y pipeline de análisis}
~\label{subsec:backend}

El backend recibe los flujos provenientes de los sensores y los 
transforma en métricas contextualizadas que alimentan el inventario 
dinámico, el motor de anomalías y la capa de exposición. Su diseño 
interno obedece al patrón arquitectónico de \textit{Pipes and 
Filters}~\cite{buschmann1996posa}: los datos atraviesan una 
secuencia de etapas especializadas conectadas por interfaces 
tipadas, y cada etapa opera sobre la salida estructurada de la 
anterior sin conocer su implementación. La Fig.~\ref{fig:servidor-central} 
ilustra esta cadena completa, desde los 
receptores TCP\@ que aceptan los flujos de los sensores hasta el 
repositorio que persiste las métricas resultantes en PostgreSQL\@.

\begin{figure}[H]
\centering
\includegraphics[width=\columnwidth]{servidor_central.png}
\caption{Pipeline interno del servidor central. Los receptores TCP aceptan los flujos de los sensores en hilos independientes; el rotador fragmenta el tráfico en archivos PCAP de ventana temporal fija; Zeek los procesa y genera logs; el parser los convierte en estructuras de conexión; el calculador aplica EWMA para derivar métricas contextualizadas, y el repositorio las persiste.}
\label{fig:servidor-central}
\end{figure}

\subsubsection{Recepción concurrente de flujos}

El servidor colector opera como punto de convergencia de datos en 
una arquitectura distribuida en la que múltiples sensores 
transmiten simultáneamente hacia el mismo destino. La concurrencia 
en el receptor es un requisito funcional: el servidor debe atender 
múltiples flujos de ingreso sin que el procesamiento de un sensor 
bloquee la recepción de datos provenientes de otros.

El modelo de concurrencia adoptado se basa en el patrón 
\textit{Thread-per-Client}~\cite{schmidt2000posa2}. Un hilo 
principal ejecuta un ciclo de aceptación de conexiones mediante la 
primitiva \texttt{accept()} del API\@ de sockets Berkeley~\cite{stevens2004unix}. 
Esta primitiva opera de forma bloqueante: 
suspende la ejecución del hilo hasta que un sensor inicia una 
conexión mediante el \textit{three-way handshake} de TCP~\cite{rfc793}. 
Al completarse el handshake, \texttt{accept()} 
retorna un descriptor de socket exclusivo para esa conexión, 
mientras el socket original permanece en estado de escucha para 
conexiones subsiguientes. Por cada conexión aceptada se instancia 
un hilo dedicado que ejecuta la lógica de lectura del protocolo 
binario descrito en la Sección~\ref{subsec:sensor}. El aislamiento 
entre hilos garantiza que la latencia o desconexión de un sensor no 
afecte la recepción de datos provenientes de los demás.

La elección de Thread-per-Client sobre alternativas como la 
multiplexación de E/S mediante \texttt{select()} o \texttt{epoll()}~\cite{stevens2004unix} 
se fundamenta en el número acotado de 
sensores previstos por los requerimientos no funcionales de 
rendimiento del sistema. Para esta escala, la simplicidad del 
modelo de hilos supera las ganancias de rendimiento que los 
mecanismos de E/S asíncronos ofrecen en escenarios de alta 
concurrencia con miles de conexiones simultáneas.

Los hilos se configuran como \textit{daemon threads}, lo que 
asegura su terminación automática cuando el proceso principal del 
servidor finaliza. Esta configuración previene la persistencia de 
hilos huérfanos que consumirían recursos del sistema operativo sin 
propósito tras un apagado controlado~\cite{silberschatz2018os}.

\subsubsection{Segmentación del tráfico en archivos PCAP}

El receptor TCP reconstruye los paquetes sanitizados emitidos por 
cada sensor y los entrega al \textit{PCAP Rotator}, el primer 
filtro del pipeline. Este componente acumula los paquetes recibidos 
en archivos con formato PCAP estándar~\cite{harris2020libpcap}, 
compuestos por un encabezado global de 24 bytes seguido de una 
secuencia de registros donde cada uno consta de un encabezado de 16 
bytes y el contenido del paquete. Este formato es el esperado por 
las herramientas convencionales de análisis de tráfico, lo que 
permite que componentes externos especializados procesen la salida 
sin necesidad de transformaciones intermedias.

La rotación obedece a una ventana temporal fija: cada archivo PCAP 
acumula el tráfico de un intervalo acotado, tras el cual se cierra 
y se emite una notificación al motor de análisis. Esta 
notificación se implementa siguiendo el patrón \textit{Observer}~\cite{gof1994}, 
que desacopla al productor del consumidor de 
eventos: el rotador publica la finalización de cada ventana sin 
necesidad de conocer qué componentes la procesarán, y el motor de 
análisis se suscribe a estas notificaciones para incorporar el 
archivo recién cerrado a su cola de trabajo. El resultado es una 
cadena donde la cadencia de análisis queda determinada por la 
ventana de rotación, no por el ritmo de llegada de los paquetes.

\subsubsection{Análisis de protocolos con Zeek}

El motor de análisis delega el procesamiento semántico del tráfico 
a Zeek, un framework de análisis de tráfico de red ampliamente 
utilizado en el ámbito de la seguridad informática~\cite{paxson1999zeek}. 
A diferencia de los analizadores orientados 
a firmas, Zeek opera sobre el comportamiento de los protocolos: 
reconstruye sesiones TCP, identifica los protocolos de aplicación 
mediante inspección de contenido, y produce registros estructurados 
que describen cada conexión en términos de sus partes, su duración, 
los bytes intercambiados en cada dirección y los metadatos 
específicos del protocolo detectado.

Cada archivo PCAP emitido por el rotador se entrega a una instancia 
de Zeek ejecutada en un contenedor Docker aislado del resto del 
sistema. Esta aislación cumple dos propósitos. El primero es de 
seguridad: el procesamiento de paquetes arbitrarios provenientes 
de una red real plantea riesgos de explotación de vulnerabilidades 
en el propio analizador, y la ejecución en contenedor acota el 
dominio de compromiso en caso de fallo. El segundo es operativo: 
Zeek y sus dependencias se mantienen versionadas de forma 
independiente del resto del backend, lo que simplifica las 
actualizaciones y garantiza la reproducibilidad del entorno de 
análisis.

Zeek produce múltiples archivos de registro por cada archivo PCAP 
procesado. Los relevantes para el pipeline son \texttt{conn.log}, 
que describe cada conexión observada con su 
tupla de direcciones y puertos, marcas temporales, duración y 
conteos de bytes y paquetes en cada sentido; \texttt{dns.log}, que 
registra las consultas y respuestas DNS junto con los tiempos de 
respuesta; y \texttt{ssl.log}, que describe los handshakes TLS 
incluyendo el nombre de servidor indicado (SNI) cuando está 
presente.

\subsubsection{Transformación de logs y cálculo de métricas}

Los archivos de registro producidos por Zeek son ficheros de texto 
tabular (TSV) que no son aptos para el análisis en memoria de 
aplicaciones posteriores. El \textit{Log Parser} actúa como 
segundo filtro del pipeline: recorre los logs, filtra las columnas 
relevantes y convierte cada línea en una estructura de datos 
tipada que representa una conexión individual. Esta estructura 
homogénea permite que el componente siguiente opere sin necesidad 
de manipular texto crudo.

El \textit{Metrics Calculator} consume las estructuras de conexión 
y deriva las métricas contextualizadas que constituyen el producto 
del pipeline. Entre ellas, el ancho de banda entrante y saliente 
por dispositivo, la tasa de retransmisiones TCP, el tiempo de 
respuesta DNS, el número de conexiones activas y fallidas, y la 
clasificación de \textit{top talkers} dentro de la ventana 
analizada. Las métricas que requieren suavizado temporal se 
computan mediante el método EWMA~\cite{lucas1990ewma,roberts1959ewma}, 
cuyo detalle formal se desarrolla en la Sección~\ref{subsec:motor} 
junto con la descripción del motor de detección de anomalías.

\subsubsection{Persistencia y separación de procesos}

Las métricas resultantes se entregan al \textit{Repository}, un 
componente que encapsula el acceso a PostgreSQL y expone una 
interfaz uniforme hacia el resto del backend. La escritura 
distingue entre entidades acumuladas y sobreescritas, según la 
tipología del modelo de datos descrita en la Sección~\ref{subsec:arquitectura}: 
cada ciclo inserta una nueva fila en 
las tablas \textsc{\_Snapshot} y actualiza mediante \texttt{UPSERT} 
las tablas \textsc{\_Current\_Metrics}.

La comunicación entre el pipeline de procesamiento y la capa de 
exposición (API REST y WebSocket, Sección~\ref{subsec:api}) se 
realiza exclusivamente a través de la base de datos, sin 
comunicación directa entre los dos procesos. Este patrón, 
conocido como \textit{Shared Database} en la literatura de 
integración empresarial~\cite{hohpe2003eip}, es apropiado cuando 
ambos extremos pertenecen a la misma organización y comparten el 
modelo de datos, y aporta dos ventajas. La primera es de 
aislamiento de fallos: un bloqueo o reinicio del pipeline no 
afecta la disponibilidad de la API para consultas móviles, y 
viceversa. La segunda es de escalabilidad independiente: ambos 
procesos pueden ajustar sus recursos asignados según su perfil de 
carga específico, que es distinto en cada caso. El pipeline opera 
en ráfagas periódicas con alto consumo de CPU, mientras que la 
API presenta un perfil de carga irregular dominado por la 
frecuencia de consulta desde el cliente móvil.

\subsection{Motor de detección de anomalías}
~\label{subsec:motor}

El motor de detección opera como consumidor de las métricas 
contextualizadas producidas por el pipeline descrito en la Sección~\ref{subsec:backend} 
y convierte cada observación en una 
evaluación cualitativa sobre su desviación respecto al comportamiento 
habitual del dispositivo. El diseño del motor se fundamenta en tres 
decisiones articuladas: un baseline adaptativo por dispositivo y por 
métrica mediante EWMA, una cuantificación de desviaciones mediante 
Z-Score adaptativo sobre varianza móvil, y un filtro de persistencia 
que exige confirmación temporal antes de generar una alerta. La 
justificación teórica de esta combinación y las limitaciones de las 
alternativas consideradas se desarrollaron en la Sección~\ref{sec:marco-teorico}; 
la presente subsección describe las 
decisiones concretas de parametrización y los detalles algorítmicos 
adoptados en la implementación.

\subsubsection{Baseline adaptativo por dispositivo}

Cada par dispositivo-métrica mantiene dos variables de estado: la 
media exponencial~\(\mu_t\) y la varianza exponencial~\(\sigma^2_t\). 
En cada ciclo de análisis (30~segundos), el motor recibe el vector 
de métricas actuales y actualiza estas dos variables mediante las 
recurrencias
\begin{equation}
\mu_t = \alpha\, x_t + (1-\alpha)\, \mu_{t-1}
\label{eq:ewma-impl}
\end{equation}
\begin{equation}
\sigma^2_t = (1-\alpha)\bigl[\sigma^2_{t-1} + \alpha(x_t - \mu_{t-1})^2\bigr]
\label{eq:ewmv-impl}
\end{equation}
con factor de suavizado~\(\alpha = 0{,}10\)~\cite{lucas1990ewma,roberts1959ewma}. 
Esta elección equilibra estabilidad y reactividad: el peso que una 
observación de hace~\(k\) ciclos retiene en el baseline decae como~\(\alpha(1-\alpha)^k\), 
de modo que una observación emitida hace 
siete ciclos (aproximadamente tres minutos y medio) conserva la mitad 
de su peso original, y una observación de hace 44 ciclos (22 minutos) 
retiene menos del uno por ciento. Más allá de esa ventana el histórico 
está efectivamente olvidado, lo que permite al motor adaptarse a 
cambios legítimos de régimen en cuestión de minutos, a diferencia del 
Z-Score con media acumulada, que requeriría horas para converger al 
nuevo nivel.

Un factor de suavizado menor (por ejemplo~\(\alpha = 0{,}01\)) 
produciría un baseline con memoria superior a cuatro horas, 
demasiado lento para adaptarse a cambios operativos legítimos de los 
dispositivos monitoreados. Un factor mayor (por ejemplo~\(\alpha = 0{,}50\)) 
produciría un baseline con memoria de tres minutos, tan 
reactivo que una anomalía sostenida se absorbería en el propio 
baseline antes de que el sistema alcanzara a reportarla con su 
magnitud completa. El valor~\(\alpha = 0{,}10\) sitúa el motor en 
un punto intermedio validado empíricamente sobre las 12~métricas 
monitoreadas por dispositivo.

\subsubsection{Varianza mínima y evaluación antes de actualización}

Dos detalles de implementación tienen consecuencias sobre el 
comportamiento del detector y merecen explicitación. El primero es 
la existencia de un piso mínimo de varianza configurado según la 
escala de cada métrica: 1{,}0 para las expresadas en porcentaje, 
100{,}0 para las expresadas en bytes por segundo, 0{,}5 para los 
conteos, 0{,}01 para los ratios y 0{,}1 para las métricas de 
latencia. Este piso se emplea en el denominador del Z-Score cuando 
la varianza acumulada cae por debajo, lo que ocurre con métricas que 
permanecen constantes durante períodos largos (por ejemplo, una tasa 
de pérdida de paquetes que se mantiene en cero durante horas) y que, 
sin ese piso, producirían Z-Scores artificialmente grandes ante 
cualquier fluctuación legítima posterior.

El segundo detalle concierne al orden de evaluación y actualización 
dentro de cada ciclo. El Z-Score se calcula contra el baseline y la 
varianza del ciclo anterior, 
\begin{equation}
Z_t = \frac{x_t - \mu_{t-1}}{\sqrt{\sigma^2_{t-1} + \varepsilon}}
\label{eq:zscore-impl}
\end{equation}
antes de actualizar esas variables con el valor observado. El 
invertir el orden haría que la anomalía se absorbiera parcialmente 
en el baseline antes de ser evaluada, reduciendo artificialmente su 
Z-Score y comprometiendo la detección. En las pruebas del sistema, 
un valor de~\(\mathtt{cpu\_pct} = 45\%\) contra un baseline de~\(15\%\) 
produce~\(Z \approx 30{,}4\) cuando se evalúa antes de 
actualizar; si se actualizaran primero el baseline y la varianza, el 
mismo valor produciría~\(Z < 3\). La evaluación anticipada preserva 
la sensibilidad completa del detector.

\subsubsection{Fase de calentamiento}

Durante los primeros 30~ciclos de observación (15~minutos) sobre un 
dispositivo nuevo, el motor acumula muestras y actualiza el baseline 
y la varianza, pero no genera alertas. La razón es cuantitativa: con 
pocas observaciones, la varianza es cercana a cero porque aún no se 
ha observado suficiente variabilidad natural. Un dispositivo cuya 
\(\mathtt{cpu\_pct}\) oscila legítimamente entre el 10~\% y el 
20~\% tendría, tras cinco ciclos, una varianza mínima; si el sexto 
valor es~22~\% (dentro del rango normal), el Z-Score explotaría a 
valores artificialmente altos porque la desviación estándar en el 
denominador sería casi nula. La ventana de 30~ciclos garantiza que 
la varianza refleje la variabilidad real del dispositivo antes de 
habilitar las alertas. Si un dispositivo se desconecta 
temporalmente y reconecta, su baseline se conserva en memoria y no 
reinicia la fase de calentamiento, dado que la desconexión no 
invalida el perfil de comportamiento previamente establecido.

\subsubsection{Clasificación de severidad}

Una vez superada la fase de calentamiento, cada ciclo produce para 
cada métrica un valor de~\(|Z_t|\) que el motor clasifica según tres 
umbrales de severidad. Las observaciones con~\(|Z_t| < 1{,}5\) se 
consideran dentro del rango normal. Aquellas con~\(1{,}5 \leq |Z_t| < 2{,}0\) 
se clasifican como informativas: se registran en la tabla 
\textsc{Insight} para análisis retrospectivo pero no generan 
notificación. Las que cumplen~\(2{,}0 \leq |Z_t| < 3{,}0\) se 
clasifican como \textit{warning}: representan desviaciones con 
probabilidad de ocurrencia de aproximadamente 4{,}56~\% bajo una 
distribución normal, una de cada 22~observaciones esperadas. Las 
que superan~\(|Z_t| \geq 3{,}0\) se clasifican como \textit{critical}: 
probabilidad de aproximadamente 0{,}27~\%, una de 
cada 370~observaciones~\cite{montgomery2013spc}. Esta gradación 
permite que cada alerta incluya una cuantificación estadística de su 
rareza, no únicamente una clasificación binaria.

\subsubsection{Filtro de persistencia}

Si cada observación con~\(|Z_t| \geq 2{,}0\) generara una alerta 
inmediata, el sistema producirían un volumen elevado de falsas 
alarmas ante fluctuaciones transitorias que son parte del 
comportamiento normal de una red LAN:@ráfagas momentáneas de 
tráfico, picos puntuales de respuesta DNS, o variaciones breves de 
latencia. Estas oscilaciones no son anomalías en el sentido 
operativo: no requieren intervención porque se resuelven por sí 
solas en el siguiente ciclo. Las anomalías reales, en contraste, se 
caracterizan por persistir durante varios ciclos consecutivos.

Para discriminar entre fluctuaciones transitorias y anomalías 
sostenidas, el motor incorpora un filtro de persistencia de 
parámetros~\(N\)~y~\(M\) en el espíritu de las reglas de Western 
Electric para cartas de control~\cite{montgomery2013spc}. El filtro 
exige que~\(|Z_t| \geq 2{,}0\) se cumpla en al menos~\(N\) de los 
últimos~\(M\) ciclos para emitir una alerta. Valores ensayados sobre 
datos sintéticos con anomalías conocidas muestran que configuraciones 
demasiado permisivas (sin filtro, o~\(N/M = 2/4\)) generan cientos 
de falsas alarmas frente a decenas de anomalías reales, mientras que 
configuraciones demasiado estrictas (\(N/M = 3/3\) o~\(3/5\)) pierden 
entre 40~\% y 45~\% de las anomalías reales.

La configuración~\(N/M = 2/2\) adoptada en el sistema exige que la 
desviación se mantenga en dos ciclos consecutivos (un minuto) antes 
de emitir la alerta. Sobre las pruebas sintéticas, esta 
configuración reduce las falsas alarmas de 383 (sin filtro) a 16, 
una disminución del 96~\%, mientras preserva 45 de las 48 anomalías 
reales inyectadas, es decir, una tasa de detección del 93{,}8~\%. 
Esta relación entre tasa de detección y tasa de falsas alarmas es la 
que justifica la elección operativa del filtro.

\subsubsection{Ciclo completo y persistencia de alertas}

El ciclo del motor se ejecuta cada 30~segundos. En cada iteración, 
para cada par dispositivo-métrica activo, el motor recupera el 
estado previo de~\(\mu_{t-1}\) y~\(\sigma^2_{t-1}\) desde memoria, 
evalúa el Z-Score contra el valor observado, actualiza el baseline 
y la varianza conforme a las ecuaciones~\eqref{eq:ewma-impl}~y~\eqref{eq:ewmv-impl}, 
incrementa el contador de ciclos consecutivos en desviación si 
aplica, y decide si emite una alerta según el filtro de persistencia. 
Las alertas confirmadas se persisten en la tabla~\textsc{Alert} 
con la severidad calculada, el valor de~\(|Z_t|\) que disparó la 
condición, la métrica implicada y la referencia al dispositivo. El 
suscriptor WebSocket descrito en la Sección~\ref{subsec:api} 
detecta la inserción y propaga la notificación al dispositivo móvil 
del administrador.

\subsection{API REST, WebSocket y red de acceso}
~\label{subsec:api}

La capa de exposición presenta los datos persistidos por el 
pipeline a la aplicación móvil del administrador. Su diseño se 
apoya en tres elementos: una interfaz REST~\cite{fielding2000rest} 
para consultas puntuales, un canal WebSocket~\cite{rfc6455} para la 
propagación de eventos en tiempo real, y una red privada Tailscale 
para el acceso remoto sin exponer el servidor a la Internet 
pública. Ambas interfaces se autentican mediante JSON Web 
Tokens~\cite{rfc7519}, lo que evita mantener estado de sesión en 
el servidor.

\subsubsection{Interfaz REST}

La API se organiza en cuatro dominios que se corresponden con las 
pantallas principales de la aplicación móvil: \texttt{/auth} para 
el login, \texttt{/network} para las métricas globales del 
segmento, \texttt{/devices} para el inventario y las métricas por 
dispositivo, y \texttt{/alerts} para las alertas generadas por el 
motor descrito en la Sección~\ref{subsec:motor}. Cada recurso se 
identifica mediante un URI semántico y se manipula con los métodos 
estándar de HTTP~\cite{rfc9110}: \texttt{GET} para consulta, 
\texttt{POST} para creación y \texttt{PATCH} para modificación 
parcial.

La implementación emplea FastAPI sobre ASGI, lo que permite 
mantener el servidor HTTP y el canal WebSocket en el mismo proceso. 
La validación de peticiones y la serialización de respuestas se 
delega a Pydantic, que verifica estáticamente los tipos declarados 
en los modelos y rechaza las peticiones malformadas antes de que 
alcancen la lógica de negocio. El contrato de la API se expone 
automáticamente como especificación OpenAPI en el endpoint 
\texttt{/docs}.

\subsubsection{Autenticación sin estado con JWT}

El login valida las credenciales contra los hashes almacenados en 
la tabla \textsc{User} y emite un token firmado con HMAC-SHA256~\cite{rfc2104}. 
El cliente lo almacena y lo presenta en el 
encabezado \texttt{Authorization} de las peticiones subsiguientes. 
La verificación se reduce a una operación criptográfica local de 
complejidad constante, por lo que no se requiere una consulta a 
PostgreSQL por cada petición. La selección de HMAC sobre firma 
asimétrica responde a que la misma entidad emite y verifica los 
tokens, escenario en el que el algoritmo simétrico es suficiente y 
más rápido.

\subsubsection{Canal WebSocket y notificaciones}

El modelo \textit{request-response} de HTTP obligaría al cliente a 
implementar \textit{polling} para detectar cambios en las métricas 
y las alertas, con el consiguiente aumento de carga sobre 
PostgreSQL y de transferencia redundante de datos. Para evitarlo, 
el sistema expone un canal WebSocket en el endpoint \texttt{/ws} 
que mantiene una conexión TCP persistente bidireccional. El cliente 
incluye el JWT como parámetro de consulta en el URI y el servidor 
valida la firma antes de aceptar el \textit{upgrade} de protocolo.

Dentro del servidor, el \textit{WebSocket Handler} se suscribe a 
los cambios de las tablas \textsc{Alert} y 
\textsc{\_Current\_Metrics} mediante el mecanismo 
\texttt{LISTEN/NOTIFY} de PostgreSQL.@Cuando el pipeline inserta o 
actualiza un registro, la base de datos emite una notificación al 
canal suscrito y el handler la propaga a los clientes conectados. 
Este flujo desacopla al pipeline de la capa de exposición: el 
pipeline desconoce a los clientes móviles y la API no sondea la 
base de datos. Cuando la aplicación móvil no está en primer plano, 
las notificaciones se entregan a través de Firebase Cloud 
Messaging, utilizando los tokens registrados en la tabla 
\textsc{Push\_Token} durante el login.

\subsubsection{Red privada con Tailscale}

El requerimiento de acceso remoto desde cualquier ubicación se 
resuelve mediante una red privada virtual construida sobre 
Tailscale, que implementa una malla entre dispositivos con el 
protocolo WireGuard~\cite{donenfeld2017wireguard}. El servidor 
central y el dispositivo móvil se registran como nodos de la misma 
red privada, reciben una dirección del rango CGNAT 
\texttt{100.64.0.0/10} y establecen un túnel cifrado directo 
cuando el estado del NAT lo permite; cuando no, el tráfico se 
enruta a través de los relays DERP de Tailscale, también cifrados. 
En ninguno de los casos el servidor expone puertos a la Internet 
pública, lo que resuelve simultáneamente la movilidad geográfica 
del administrador y la confidencialidad del canal de comunicación 
sin degradar la postura de seguridad del backend.

\subsection{Integración continua, pruebas y despliegue}
~\label{subsec:cicd}

La exposición de los componentes descritos en las secciones
anteriores se sostiene sobre un conjunto de pipelines automatizados
que validan cada incremento, producen los artefactos distribuibles
y los publican en los registros correspondientes. Esta subsección
describe las decisiones técnicas que materializan la práctica
metodológica introducida en la Sección~\ref{sec:metodologia}: la
estructura de los cuatro workflows alojados en los dos repositorios,
la batería de pruebas ejecutada por cada uno, y la infraestructura
de despliegue contenedorizado que recibe los artefactos resultantes.

\subsubsection{Visión global}

Cada uno de los dos repositorios mantiene dos workflows de GitHub
Actions: uno de integración continua (CI) que se dispara sobre todo
\texttt{push} y toda \textit{pull request} hacia \texttt{main}, y
uno de entrega continua (CD) que se dispara únicamente sobre
\texttt{push} a \texttt{main}. Los cuatro workflows imponen un
límite temporal de diez minutos (\texttt{timeout-minutes: 10}) que
acota el costo de una ejecución fallida y preserva el principio de
retroalimentación rápida. Los workflows de la aplicación móvil
incorporan además directivas de concurrencia que cancelan
ejecuciones obsoletas cuando se registra un nuevo \texttt{push}
sobre la misma rama, y filtros de rutas que evitan disparar el
pipeline ante cambios exclusivamente documentales.

La separación CI/CD responde a una decisión arquitectónica clásica
documentada por Humble y Farley~\cite{humble2010cd}: validar sin
producir artefactos acelera el ciclo de revisión, y producir
artefactos sólo sobre la rama aceptada evita la proliferación de
imágenes o APKs asociados a ramas de trabajo que eventualmente
serán descartadas. La reproducibilidad se preserva mediante tres
mecanismos: las versiones de las acciones se fijan explícitamente
en los archivos de workflow, las dependencias se cachean con claves
derivadas del hash del archivo de lock correspondiente, y las
credenciales sensibles se inyectan como secretos del repositorio
sin aparecer nunca en el código.

\subsubsection{Pipeline del backend}

El workflow \texttt{ci.yml} del backend ejecuta sobre una imagen
\texttt{ubuntu-latest} cuatro etapas secuenciales. Tras el checkout
del repositorio, instala Python 3.12 y restaura la caché de
\texttt{pip} con clave derivada del hash de
\texttt{requirements*.txt}. A continuación instala las dependencias
de desarrollo, ejecuta el linter \texttt{ruff} sobre los directorios
\texttt{src/} y \texttt{tests/}, y corre la suite completa con
\texttt{pytest tests/ -v -{}-cov=src -{}-cov-report=xml
-{}-cov-report=term-missing}, exportando \texttt{PYTHONPATH=.} como
variable de entorno para que los tests puedan importar el paquete
\texttt{src}. El reporte de cobertura se publica como artefacto del
workflow con la condición \texttt{if: always()}, que asegura su
disponibilidad incluso cuando la fase de pruebas falla.

El workflow \texttt{cd.yml} se ejecuta únicamente sobre
\texttt{push} a la rama principal y asume que CI ha concluido
satisfactoriamente. Sus permisos declarativos
(\texttt{packages: write}) habilitan la publicación en el registro
de contenedores de GitHub. Tras autenticarse en \texttt{ghcr.io}
con el \texttt{GITHUB\_TOKEN} inyectado automáticamente por la
plataforma, invoca la acción \texttt{docker/build-push-action}
sobre el \texttt{Dockerfile} de la raíz del repositorio y publica
dos etiquetas simultáneas: el SHA del commit para trazabilidad
exacta y \texttt{latest} para conveniencia del consumidor. La
elección del GitHub Container Registry sobre alternativas como
Docker Hub responde a dos criterios prácticos: la autenticación se
resuelve con el token implícito del repositorio sin configurar
secretos adicionales, y los límites de descarga resultan
considerablemente más laxos en entornos automatizados donde cada
ejecución del pipeline podría traducirse en una descarga.

El \texttt{Dockerfile} del backend adopta una estrategia de
construcción multi-etapa. Una primera etapa basada en
\texttt{python:3.12-slim} instala las dependencias en un prefijo
aislado; una segunda etapa, también basada en la imagen slim, copia
exclusivamente el directorio de dependencias ya compiladas y el
código fuente. Este patrón reduce el tamaño final de la imagen al
descartar herramientas de compilación y cachés intermedias, y
define además un \texttt{HEALTHCHECK} que ejercita el endpoint
\texttt{/docs} de FastAPI periódicamente, habilitando a los
orquestadores a detectar instancias no saludables y reemplazarlas
automáticamente.

\subsubsection{Pipeline de la aplicación móvil}

El workflow \texttt{ci.yml} de la aplicación móvil replica la
estructura del backend adaptada al stack Flutter. Instala el SDK
mediante la acción \texttt{subosito/flutter-action} con canal
\texttt{stable} y caché habilitada, restaura la caché de paquetes
\texttt{pub} con clave derivada de \texttt{pubspec.lock}, y ejecuta
tres validaciones sucesivas. El análisis estático se invoca con
\texttt{flutter analyze -{}-no-fatal-infos}, configuración que
habilita la visibilidad de advertencias informativas sin
convertirlas en fallos del pipeline durante la fase de
estabilización del proyecto. La verificación de formato emplea
\texttt{dart format -{}-output=none -{}-set-exit-if-changed}, que
aborta el workflow ante cualquier desviación respecto al formato
canónico del lenguaje, preservando así una invariante de estilo
sin intervención manual. Finalmente, \texttt{flutter test
-{}-coverage} ejecuta las suites unitarias y de widget generando un
reporte \texttt{lcov} que se publica como artefacto.

El workflow \texttt{cd.yml} construye un APK de depuración sobre
cada merge aceptado. Antes de invocar \texttt{flutter build}, el
workflow materializa el archivo \texttt{google-services.json}
requerido por Firebase: si el secreto
\texttt{GOOGLE\_SERVICES\_JSON} está definido en el repositorio, su
contenido codificado en base64 se decodifica y escribe en
\texttt{android/app/google-services.json}; en caso contrario, el
workflow genera un documento placeholder que permite completar la
compilación aunque las notificaciones push no funcionen en el
artefacto resultante. La URL del backend se inyecta como constante
de compilación mediante la bandera \texttt{-{}-dart-define}, que
vincula la constante \texttt{API\_BASE\_URL} con el valor del
secreto homónimo del repositorio, permitiendo que el mismo binario
apunte a distintos entornos sin modificar el código fuente. El APK resultante se publica como
artefacto del workflow con retención de catorce días y nombre
derivado del SHA del commit, preservando la trazabilidad hacia el
cambio que lo originó.

La decisión de empaquetar en modo depuración obedece a las
restricciones del proyecto académico: no existe firma de
publicación configurada y el artefacto se distribuye exclusivamente
para demostración y pruebas internas. Un despliegue en producción
exigiría reemplazar este paso por \texttt{flutter build apk
-{}-release} con una configuración de firma que consuma un
\texttt{keystore} inyectado como secreto adicional, tal como se
reconoce en la Sección~\ref{sec:discusion} entre los elementos que
separan el prototipo académico del producto operable.

\subsubsection{Estrategia de pruebas}

La pirámide de pruebas del proyecto se construye con profundidades
distintas en cada componente, reflejando la relevancia relativa de
la lógica de negocio frente a la capa de presentación. La
Tabla~\ref{tab:suites-pruebas} sintetiza las suites ejecutadas por
los pipelines de CI y su distribución por dominio funcional.

\begin{table}[ht]
\centering
\caption{Distribución de suites de pruebas por componente y dominio.}
\label{tab:suites-pruebas}
\footnotesize
\begin{tabular}{@{}p{1.8cm}p{3.5cm}c@{}}
\toprule
\textbf{Componente} & \textbf{Dominio} & \textbf{Suites} \\
\midrule
Backend & Autenticación, 2FA, scopes, recovery & 12 \\[2pt]
Backend & Motor de anomalías y pipeline & 4 \\[2pt]
Backend & Endpoints, rate limiting, rendimiento & 4 \\[2pt]
Backend & Seguridad y configuración & 3 \\[2pt]
Backend & Infraestructura (factory, migrations, edge) & 4 \\[2pt]
App móvil & Providers y servicios & 2 \\[2pt]
App móvil & Pantallas TOTP y registro & 3 \\[2pt]
App móvil & Widget smoke del árbol de inyección & 1 \\
\bottomrule
\end{tabular}
\end{table}

En el backend, el aislamiento entre pruebas se consigue mediante la
sobreescritura de dependencias de FastAPI con
\texttt{app.dependency\_overrides} y una fixture de sesión por test
que construye y destruye el esquema de base de datos en cada
ejecución. Esta fixture opera por defecto sobre SQLite en memoria,
lo que permite ciclos de ejecución del orden de decenas de segundos
sobre la suite completa. Una única suite,
\texttt{test\_auth\_race\_postgres.py}, constituye la excepción:
ejerce un escenario de concurrencia en el consumo de tokens de
invitación que depende de la semántica precisa de \texttt{SELECT
\ldots FOR UPDATE}~\cite{silberschatz2018os} y por tanto requiere
una instancia real de PostgreSQL. Esta separación entre pruebas
rápidas sobre SQLite y pruebas puntuales sobre el motor real
ilustra una decisión deliberada: la velocidad del ciclo de
desarrollo se prioriza en el caso general, y el costo adicional de
levantar el motor real se asume únicamente en el punto donde la
semántica de SQLite difiere de la de producción.

En la aplicación móvil, las pruebas se concentran en la máquina de
estados del proveedor de autenticación y en el servicio HTTP. La
fidelidad al comportamiento real se preserva mockeando
exclusivamente las fronteras del sistema ---el cliente \texttt{Dio}
y el canal \texttt{MethodChannel} que intermedia con el
\texttt{Keystore} del dispositivo--- mientras que la lógica de
dominio se ejerce sin sustitutos. Un único test de widget,
\texttt{widget\_test.dart}, verifica que el árbol completo de
inyección de dependencias se instancie sin errores al arranque,
cumpliendo la función de prueba de integración de cableado aunque
no de interacción con la interfaz.

La cobertura instrumentada se reporta como artefacto del pipeline
pero no se utiliza como criterio bloqueante. Esta decisión responde
al contexto académico del proyecto, donde imponer un umbral rígido
durante la fase de estabilización introduciría fricción sin elevar
la calidad efectiva del código. Los objetivos declarados para la
revisión sobre los reportes generados al cierre de cada iteración
se sitúan en aproximadamente el setenta y cinco por ciento de
líneas cubiertas en el módulo \texttt{src/} del backend y del
sesenta por ciento en los módulos \texttt{lib/core/} y
\texttt{lib/infrastructure/} de la aplicación móvil.

\subsubsection{Dashboard de validación del motor de anomalías}

Además de la suite \texttt{pytest} ejecutada en CI, el repositorio
del backend incluye el script \texttt{tests/generate\_report.py},
concebido como herramienta de evidencia auditable para la
sustentación del proyecto. El script ejecuta los mismos
componentes del motor de detección descritos en la
Sección~\ref{subsec:motor} ---el baseline adaptativo, el detector
por dispositivo y el generador de insights--- sobre series
sintéticas con eventos conocidos: calentamiento sobre treinta
ciclos, inyección de \textit{spikes} en el consumo de CPU,
verificación de la conversión entre Z-Score y probabilidad, y
validación del ciclo completo de autenticación con emisión y
verificación de JWT.

El resultado se serializa como un documento HTML autocontenido,
\texttt{validation\_report.html}, que presenta cada ciclo con su
valor observado, el baseline estimado, la varianza y el Z-Score
resultante, junto con la clasificación de severidad aplicada por
el motor. Este artefacto complementa la cobertura numérica
reportada por CI con una visualización de la lógica algorítmica
sobre datos conocidos, y permite a un revisor externo inspeccionar
el comportamiento del detector sin necesidad de ejecutar el
pipeline completo ni acceder al tráfico de red real.

\subsubsection{Despliegue contenedorizado}

El despliegue del backend se apoya en la imagen OCI publicada por
\texttt{cd.yml} sobre \texttt{ghcr.io}. El archivo
\texttt{docker-compose.yml} presente en el repositorio del servidor
orquesta un servicio de PostgreSQL 15 con \textit{healthcheck} y
volumen persistente \texttt{pgdata}, que constituye la dependencia
compartida entre los dos procesos independientes descritos en la
Sección~\ref{subsec:backend}. El motor de pipeline y la API REST
no se incluyen en el mismo \texttt{docker-compose.yml} y se
despliegan como contenedores separados a partir de la imagen
publicada, coherente con la decisión arquitectónica de aislar sus
ciclos de vida y sus perfiles de consumo de recursos.

Un consumidor autorizado de la imagen puede desplegar el backend
en cualquier entorno compatible con el runtime OCI mediante una
secuencia de \texttt{docker pull} sobre la etiqueta del commit
correspondiente, seguida de \texttt{docker run} con las variables
de entorno \texttt{DATABASE\_URL} y \texttt{JWT\_SECRET} inyectadas
en tiempo de arranque y con el puerto 8000 expuesto. Estas
variables desacoplan la imagen de cualquier valor específico de
entorno y permiten promover la misma imagen probada en CI hacia
distintos entornos de despliegue sin recompilación. Esta
propiedad, conocida en la literatura de entrega continua como
\textit{build once, deploy anywhere}~\cite{humble2010cd}, cierra
el ciclo entre el cambio aceptado en el repositorio y la instancia
en ejecución sobre la que se valida el comportamiento del sistema.